\documentclass[11pt]{article}
\usepackage{amsmath,amssymb,amsthm,bm}
\usepackage{geometry}
\geometry{margin=1in}
\usepackage{setspace}
\usepackage{enumitem}
\usepackage{booktabs}
\usepackage{mathtools}
\usepackage{hyperref}
\usepackage{algorithm}
\usepackage{algpseudocode}

\title{Joint Multilayered LSIRM with Probabilistic-PCA Mixture Clustering\\
NIW Predictive Allocation and Split--Merge Moves\\
(Model V10)}
\date{May 1, 2026}

\begin{document}
\maketitle

\section{Overview and motivation}

Model V10 extends Model V9 by modifying the clustering layer of the
PPCA-mixture block.  The PPCA measurement equation introduced in V9 is
retained:
\[
    x_j \mid \delta,\Lambda,\eta_j,\sigma_\varepsilon^2
    \sim N_n(\delta+\Lambda\eta_j,\sigma_\varepsilon^2 I_n),
\]
where $x_j$ is the row-centered LSIRM-induced log-distance profile of item
$j$.  The V9 change from per-respondent residual variances
$\psi_1,\dots,\psi_n$ to a single scalar $\sigma_\varepsilon^2$ is kept,
because it prevents the FA layer from absorbing the entire respondent-wise
residual structure and thereby helps maintain a non-degenerate loading
matrix $\Lambda$.

The new change in V10 is not in the PPCA measurement equation, but in the
mixture clustering step for the item factor scores $\eta_j$.  V9 used a
standard conditional Gaussian-mixture update
\[
    P(c_j=l\mid\cdot)\propto \rho_l\,\phi_r(\eta_j;\mu_l,\Sigma_l),
\]
with independent priors $\mu_l\sim N_r(m_0,V_0)$ and
$\Sigma_l\sim IW(\nu_0,S_0)$.  This update is local and can be sticky: once
an overfitted component becomes empty, the sampled value of
$(\mu_l,\Sigma_l)$ is often poorly located, and the component may never be
reoccupied.  In addition, single-site Gibbs updates move only one item at a
time and can have difficulty splitting a large cluster into two coherent
subclusters.

V10 addresses this with two changes:
\begin{enumerate}[leftmargin=2em]
    \item Replace the independent $(\mu_l,\Sigma_l)$ prior by a conjugate
    Normal--Inverse-Wishart (NIW) prior, and update $c_j$ using the NIW
    posterior predictive density, integrating out $(\mu_l,\Sigma_l)$ and
    $\rho$ for the purpose of cluster allocation.
    \item Add split--merge Metropolis--Hastings moves on the partition
    $c=(c_1,\dots,c_P)$, using the same NIW-collapsed target.  These moves
    propose major changes to the partition by splitting one occupied
    component into two or merging two occupied components into one.
\end{enumerate}

The purpose is to reduce label stickiness and improve exploration of
partitions with $K_+>2$, while keeping the V9 PPCA identification
improvement.

\section{Notation}

\subsection*{Carried over from V9}

\begin{itemize}[leftmargin=1.5em]
    \item $n,L,P_\ell,P,d,r,K^\star$: number of respondents, layers, items,
    LSIRM dimension, factor dimension, and maximum number of mixture
    components.
    \item $a_i^{(\ell)},b_j^{(\ell)}$: LSIRM respondent and item positions.
    \item $x_{ij}=\log\|a_i^{(\ell(j))}-b_j^{(\ell(j))}\|$: LSIRM-induced
    log-distance profile, row-centered at each MCMC iteration.
    \item $x_j=(x_{1j},\dots,x_{nj})^\top\in\mathbb R^n$: item $j$'s
    row-centered profile.
    \item $\eta_j\in\mathbb R^r$: PPCA factor score for item $j$.
    \item $c_j\in\{1,\dots,K^\star\}$: cluster label of item $j$.
    \item $\Lambda\in\mathbb R^{n\times r}$: loading matrix; row $i$ is
    $\lambda_i^\top$.
    \item $\delta=(\delta_1,\dots,\delta_n)^\top$: respondent baseline.
    \item $\sigma_\varepsilon^2$: shared PPCA residual variance.
    \item $\rho\in\Delta^{K^\star-1}$: finite mixture weights.
    \item $\mu_l\in\mathbb R^r$, $\Sigma_l\in\mathbb R^{r\times r}_{\succ0}$:
    component-specific mean and covariance of the $\eta$ mixture.
\end{itemize}

\subsection*{New or modified in V10}

\begin{itemize}[leftmargin=1.5em]
    \item $\kappa_0>0$: NIW prior precision for the component mean.
    Replaces the separate Gaussian prior covariance $V_0$ used in V9.
    \item For a cluster $l$, define
    \[
        \mathcal J_l=\{j:c_j=l\},\qquad n_l=|\mathcal J_l|.
    \]
    \item For leave-one-out calculations,
    \[
        \mathcal J_{l,-j}=\{m:m\neq j,\ c_m=l\},\qquad
        n_{l,-j}=|\mathcal J_{l,-j}|.
    \]
    \item $m_l,\kappa_l,\nu_l,S_l$: NIW posterior parameters for component
    $l$.
    \item $m_{l,-j},\kappa_{l,-j},\nu_{l,-j},S_{l,-j}$: leave-one-out NIW
    posterior parameters used for updating $c_j$.
\end{itemize}

\subsection*{Removed or no longer central}

\begin{itemize}[leftmargin=1.5em]
    \item $V_0$ is removed from the mixture prior and replaced by
    $\kappa_0$ through
    $\mu_l\mid\Sigma_l\sim N_r(m_0,\Sigma_l/\kappa_0)$.
    \item $\rho$ is no longer used directly in the single-site update of
    $c_j$.  It is integrated out, yielding the factor $n_{l,-j}+e_0$.
    A draw of $\rho$ may still be sampled after label updates for monitoring
    and storage.
\end{itemize}

\section{LSIRM block}

The LSIRM block is unchanged from V9.  The LSIRM parameters are updated
without including the PPCA-mixture likelihood in the Metropolis--Hastings
acceptance ratios.  Thus, the coupling remains one-way:
\[
    \text{LSIRM positions}\quad\longrightarrow\quad
    X=(x_{ij})\quad\longrightarrow\quad
    \text{PPCA-mixture clustering}.
\]
At each MCMC iteration, the current LSIRM positions generate
\[
    x_{ij}=\log\max\{\|a_i^{(\ell(j))}-b_j^{(\ell(j))}\|,\epsilon\},
\]
and row-centering is applied before the PPCA-mixture block is updated.

\section{PPCA-mixture clustering block}

Conditional on the current LSIRM state, the V10 measurement equation is
\begin{equation}
    x_j\mid\delta,\Lambda,\eta_j,\sigma_\varepsilon^2
    \sim
    N_n(\delta+\Lambda\eta_j,\sigma_\varepsilon^2 I_n),
    \qquad j=1,\dots,P.
    \label{eq:v10-measurement}
\end{equation}

The factor scores are assigned a finite Gaussian mixture prior:
\begin{align}
    \eta_j\mid c_j=l,\mu_l,\Sigma_l
    &\sim N_r(\mu_l,\Sigma_l),\label{eq:v10-eta-mixture}\\
    c_j\mid\rho
    &\sim \mathrm{Categorical}(\rho_1,\dots,\rho_{K^\star}),\\
    \rho
    &\sim \mathrm{Dirichlet}(e_0,\dots,e_0).
\end{align}

The component parameters now follow an NIW prior:
\begin{align}
    \Sigma_l &\sim IW(\nu_0,S_0),\label{eq:niw-sigma-prior}\\
    \mu_l\mid\Sigma_l &\sim
    N_r\left(m_0,\frac{1}{\kappa_0}\Sigma_l\right),
    \qquad l=1,\dots,K^\star.\label{eq:niw-mu-prior}
\end{align}

\section{Priors}

The non-mixture priors remain as in V9:
\begin{align}
    \delta_i\mid\sigma_\delta^2 &\sim N(0,\sigma_\delta^2), &
    \sigma_\delta^2 &\sim IG(a_\delta,b_\delta),\\
    \lambda_i &\sim N_r(0,\tau_\lambda^2 I_r),\\
    \sigma_\varepsilon^2 &\sim IG(a_\varepsilon,b_\varepsilon).
\end{align}

The mixture prior is changed from independent $N$ and $IW$ priors to the
conjugate NIW prior in \eqref{eq:niw-sigma-prior}--\eqref{eq:niw-mu-prior}.
A practical default set is:
\begin{center}
\begin{tabular}{lll}
\toprule
Parameter & Default & Role \\
\midrule
$e_0$ & $0.1$ & Sparse overfitted finite mixture prior. \\
$m_0$ & $0_r$ & Prior center of component means. \\
$\kappa_0$ & $0.001$ & Weak prior precision for $\mu_l$ (matches $V_0\!\approx\!9I_r$ in V9). \\
$\nu_0$ & $r+10$ & Inverse-Wishart degrees of freedom. \\
$S_0$ & $0.01 I_r$ & Prior within-cluster covariance scale (see \S\ref{sec:s0-elicitation}). \\
$\tau_\lambda^2$ & $0.25$ & Loading prior variance. \\
$(a_\delta,b_\delta)$ & $(2,1)$ & Prior for respondent baselines. \\
$(a_\varepsilon,b_\varepsilon)$ & $(5,0.1)$ & Prior for shared PPCA residual variance. \\
\bottomrule
\end{tabular}
\end{center}

The choice of $\kappa_0$ should be small if $m_0=0$ is meant to be weakly
informative.  The previous V9 prior $\mu_l\sim N_r(m_0,V_0)$ can be matched
approximately by setting $\kappa_0$ so that the marginal magnitude of
$\mu_l$ matches $V_0$:
\[
    E[\Sigma_l/\kappa_0]
    =\frac{S_0}{(\nu_0-r-1)\kappa_0}
    \;\stackrel{!}{=}\;V_0.
\]
With $\nu_0=r+10$, $S_0=0.01I_r$ and $V_0=9I_r$, this gives
$\kappa_0=0.01/(9\cdot 9)\approx 1.2\times 10^{-4}$; a default of
$\kappa_0=10^{-3}$ keeps $\mu_l$ weakly informative while remaining within
an order of magnitude of the matched value.

\subsection{Eliciting $S_0$: matching the within-cluster scale of $\eta$}
\label{sec:s0-elicitation}

The default $S_0$ value above is the most consequential of the NIW
hyperparameters.  In the Lamb-style finite-mixture family, a too-loose
$S_0$ is the dominant failure mode: it makes the model believe components
are wider than they actually are, inflates the NIW marginal likelihood for
"loose" clusters, and suppresses splits even when the factor scores
clearly support multiple components.  The diagnostic and remedy are:

\begin{enumerate}[leftmargin=2em]
    \item \textbf{Run a short pilot chain.}  After warm-up, save the
    posterior draws of $\eta_j$ (or compute the posterior mean
    $\bar\eta_j$).  Compute the empirical within-cluster sd along each
    factor dimension using the pilot's modal $c_j$:
    \[
        \widehat{\mathrm{sd}}_q
        =\sqrt{\tfrac{1}{P-K_+}\sum_l\sum_{j\in\mathcal J_l}
        (\eta_{jq}-\bar\eta_{lq})^2},
        \qquad q=1,\dots,r.
    \]
    \item \textbf{Compare to the prior implied scale.}  Under
    $\Sigma_l\sim IW(\nu_0,S_0)$ with $S_0=s_0 I_r$,
    \[
        E[\Sigma_l]=\frac{s_0}{\nu_0-r-1}\,I_r,
        \qquad
        \text{prior implied per-dim sd}\approx\sqrt{\tfrac{s_0}{\nu_0-r-1}}.
    \]
    \item \textbf{Tighten if the prior dominates.}  A factor of $5$--$10$
    over the empirical sd is a flag: the prior expects clusters to be
    that much wider than the data suggests, so split moves will be
    suppressed.  Reset $s_0$ so that the prior implied per-dim sd is on
    the same order as $\widehat{\mathrm{sd}}_q$.
\end{enumerate}

\paragraph{Concrete example.}  In a four-layer simulation
($n=150$, $P=40$, $K_{\rm true}=4$ diamond clusters), a pilot chain at
$r=2$, $S_0=0.05I_r$, $\nu_0=r+10=12$ gave empirical
$\widehat{\mathrm{sd}}_q\approx 0.01$--$0.02$; the prior implied per-dim
sd was $\sqrt{0.05/9}\approx 0.075$, i.e.\ a $4$--$7\times$ overestimate.
Splits were rare (split acceptance $\approx 10\%$) and the chain stalled
at $K_+=2$--$3$ with $\mathrm{ARI}\approx 0.46$.  Refitting with
$S_0=0.01I_r$ (prior implied per-dim sd $\approx 0.033$) gave
$K_+=5$, split acceptance $15\%$, and $\mathrm{ARI}=0.89$.  The default
$S_0=0.01I_r$ in the table above reflects this calibration; users on
data with a different $\eta$ scale should re-run this elicitation rather
than treat $0.01I_r$ as universal.

\section{Joint hierarchical model}

The full joint posterior is proportional to
\begin{align*}
&p(\Theta_{\mathrm{LSIRM}},\Theta_{\mathrm{clust}}\mid Y)\\
&\quad\propto
p(Y\mid\Theta_{\mathrm{LSIRM}})\\
&\qquad\times
\prod_{j=1}^P
N_n\{x_j(\Theta_{\mathrm{LSIRM}});\delta+\Lambda\eta_j,
\sigma_\varepsilon^2 I_n\}
N_r(\eta_j;\mu_{c_j},\Sigma_{c_j})\rho_{c_j}\\
&\qquad\times
p(\Theta_{\mathrm{LSIRM}})
\,p(\delta\mid\sigma_\delta^2)p(\sigma_\delta^2)
\,p(\Lambda)p(\sigma_\varepsilon^2)
\,p(\rho)
\prod_{l=1}^{K^\star}p(\mu_l,\Sigma_l),
\end{align*}
where
\[
    p(\mu_l,\Sigma_l)
    =p(\mu_l\mid\Sigma_l)p(\Sigma_l).
\]
This enforces the NIW prior.  The clustering parameter set is
\[
\Theta_{\mathrm{clust}}
=\{\eta_j,c_j,\rho,\delta,\Lambda,\sigma_\varepsilon^2,
\mu_l,\Sigma_l,\sigma_\delta^2\}.
\]

\section{NIW posterior quantities}

For a generic index set $A\subseteq\{1,\dots,P\}$, let
\[
    n_A=|A|,
    \qquad
    \bar\eta_A=\frac{1}{n_A}\sum_{j\in A}\eta_j
    \quad(n_A>0),
\]
and
\[
    Q_A=\sum_{j\in A}(\eta_j-\bar\eta_A)(\eta_j-\bar\eta_A)^\top
    \quad(n_A>0).
\]
For $n_A=0$, define $Q_A=0$ and use the prior parameters directly.

The NIW posterior parameters for cluster $A$ are
\begin{align}
    \kappa_A &= \kappa_0+n_A,\\
    \nu_A &= \nu_0+n_A,\\
    m_A &=
    \begin{cases}
    \dfrac{\kappa_0m_0+n_A\bar\eta_A}{\kappa_0+n_A}, & n_A>0,\\[1.25em]
    m_0, & n_A=0,
    \end{cases}\\
    S_A &=
    \begin{cases}
    S_0+Q_A+
    \dfrac{\kappa_0 n_A}{\kappa_0+n_A}
    (\bar\eta_A-m_0)(\bar\eta_A-m_0)^\top, & n_A>0,\\[1.25em]
    S_0, & n_A=0.
    \end{cases}
\end{align}

For cluster $l$, write these quantities as
\[
    (\kappa_l,\nu_l,m_l,S_l)
    = (\kappa_{\mathcal J_l},\nu_{\mathcal J_l},m_{\mathcal J_l},S_{\mathcal J_l}).
\]
For leave-one-out allocation of item $j$, write
\[
    (\kappa_{l,-j},\nu_{l,-j},m_{l,-j},S_{l,-j})
    = (\kappa_{\mathcal J_{l,-j}},\nu_{\mathcal J_{l,-j}},
       m_{\mathcal J_{l,-j}},S_{\mathcal J_{l,-j}}).
\]

\section{NIW posterior predictive density}

The leave-one-out posterior predictive distribution of $\eta_j$ under
component $l$ is multivariate Student-$t$:
\begin{equation}
    \eta_j\mid \eta_{\mathcal J_{l,-j}}
    \sim
    t_{d_{l,-j}}
    \left(
    m_{l,-j},
    V^{\mathrm{pred}}_{l,-j}
    \right),
    \label{eq:niw-predictive}
\end{equation}
where
\begin{align}
    d_{l,-j} &= \nu_{l,-j}-r+1,\\
    V^{\mathrm{pred}}_{l,-j}
    &=
    \frac{\kappa_{l,-j}+1}
    {\kappa_{l,-j}(\nu_{l,-j}-r+1)}S_{l,-j}.
\end{align}

Thus the collapsed single-site allocation update is
\begin{equation}
    P(c_j=l\mid c_{-j},\eta)
    \propto
    (n_{l,-j}+e_0)
    t_{d_{l,-j}}
    \left(
    \eta_j;
    m_{l,-j},
    V^{\mathrm{pred}}_{l,-j}
    \right),
    \qquad l=1,\dots,K^\star.
    \label{eq:collapsed-c-update}
\end{equation}

For an empty component, $n_{l,-j}=0$, so
\[
    (\kappa_{l,-j},\nu_{l,-j},m_{l,-j},S_{l,-j})
    =(\kappa_0,\nu_0,m_0,S_0),
\]
and the allocation probability is based on the prior predictive density
multiplied by $e_0$.  This is the mechanism by which empty components can be
reoccupied.

\section{Collapsed partition target for split--merge moves}

Let $c=(c_1,\dots,c_P)$ be the current label vector and
$\mathcal J_l=\{j:c_j=l\}$.  Integrating out $\rho$ and
$(\mu_l,\Sigma_l)_{l=1}^{K^\star}$ gives the collapsed target
\begin{equation}
    \pi(c\mid\eta)
    \propto
    \prod_{l=1}^{K^\star}
    \frac{\Gamma(n_l+e_0)}{\Gamma(e_0)}
    \;m(\eta_{\mathcal J_l}),
    \label{eq:collapsed-target}
\end{equation}
where
\[
    m(\eta_{\mathcal J_l})
    =
    \int
    \prod_{j\in\mathcal J_l}N_r(\eta_j;\mu_l,\Sigma_l)
    p(\mu_l,\Sigma_l)\,d\mu_l\,d\Sigma_l
\]
is the NIW marginal likelihood of the items in cluster $l$.

For $A\subseteq\{1,\dots,P\}$, the log NIW marginal likelihood is
\begin{align}
    \log m(\eta_A)
    &=-\frac{n_A r}{2}\log\pi
    +\frac{r}{2}\{\log\kappa_0-\log\kappa_A\}
    +\frac{\nu_0}{2}\log|S_0|
    -\frac{\nu_A}{2}\log|S_A| \notag\\
    &\qquad
    +\log\Gamma_r\left(\frac{\nu_A}{2}\right)
    -\log\Gamma_r\left(\frac{\nu_0}{2}\right),
    \label{eq:niw-marginal}
\end{align}
The factor $\pi^{-n_A r/2}$ (rather than $(2\pi)^{-n_A r/2}$) follows from
the cancellation between the Gaussian likelihood normalizing constant and
the inverse-Wishart normalizing constant; see Murphy (2007), eq.~(232).
where
\[
    \log\Gamma_r(a)
    =\frac{r(r-1)}{4}\log\pi
    +\sum_{q=1}^r\log\Gamma\left(a+\frac{1-q}{2}\right).
\]

Up to constants independent of $c$, the log collapsed partition target is
\begin{equation}
    \log\pi(c\mid\eta)
    =
    \sum_{l=1}^{K^\star}
    \left[
    \log\Gamma(n_l+e_0)-\log\Gamma(e_0)
    +\log m(\eta_{\mathcal J_l})
    \right]
    +\mathrm{constant}.
    \label{eq:log-collapsed-target}
\end{equation}

\section{Posterior updates}

LSIRM updates are unchanged.  The PPCA-mixture updates are modified as
follows.

\paragraph{(1) $\eta_j$.}
Given the current component parameters, the conditional posterior remains
Gaussian:
\begin{equation}
    \eta_j\mid\cdot\sim N_r(m_{\eta_j},V_{\eta_j}),
\end{equation}
where
\begin{align}
    V_{\eta_j}
    &=
    \left(\sigma_\varepsilon^{-2}\Lambda^\top\Lambda+
    \Sigma_{c_j}^{-1}\right)^{-1},\\
    m_{\eta_j}
    &=
    V_{\eta_j}
    \left\{
    \sigma_\varepsilon^{-2}\Lambda^\top(x_j-\delta)
    +\Sigma_{c_j}^{-1}\mu_{c_j}
    \right\}.
    \label{eq:eta-update-v10}
\end{align}

\paragraph{(2) $c_j$: NIW-collapsed single-site update.}
Remove item $j$ from its current cluster and compute the leave-one-out NIW
posterior parameters for every component $l$.  Then sample
\begin{equation}
    P(c_j=l\mid c_{-j},\eta)
    \propto
    (n_{l,-j}+e_0)
    t_{\nu_{l,-j}-r+1}
    \left(
    \eta_j;
    m_{l,-j},
    \frac{\kappa_{l,-j}+1}
    {\kappa_{l,-j}(\nu_{l,-j}-r+1)}S_{l,-j}
    \right).
    \label{eq:c-update-v10}
\end{equation}
This replaces the V9 update
$P(c_j=l\mid\cdot)\propto \rho_l\phi_r(\eta_j;\mu_l,\Sigma_l)$.

\paragraph{(3) Split--merge move for $c$.}
After a full sweep of single-site NIW-collapsed updates, optionally perform
$M_{\mathrm{SM}}$ split--merge proposals targeting the collapsed density
\eqref{eq:log-collapsed-target}.  Details are given in
Section~\ref{sec:split-merge}.

\paragraph{(4) $\Sigma_l$ and $\mu_l$: NIW posterior draw.}
Given the final labels after the single-site and split--merge updates,
compute $(\kappa_l,\nu_l,m_l,S_l)$ for each component.  Then sample
\begin{align}
    \Sigma_l\mid\eta,c
    &\sim IW(\nu_l,S_l),\\
    \mu_l\mid\Sigma_l,
    \eta,c
    &\sim N_r\left(m_l,\frac{1}{\kappa_l}\Sigma_l\right).
    \label{eq:mu-sigma-update-v10}
\end{align}
For empty components, this reduces to a draw from the NIW prior.

\paragraph{(5) $\rho$ for storage or monitoring.}
Although $\rho$ is integrated out in the label update, it may be sampled
after the labels are updated:
\begin{equation}
    \rho\mid c
    \sim
    \mathrm{Dirichlet}(e_0+n_1,\dots,e_0+n_{K^\star}).
    \label{eq:rho-storage-v10}
\end{equation}
This draw is not needed for $c_j$ updates, but is useful for posterior
summaries and compatibility with existing output structure.

\paragraph{(6) $\delta_i$.}
The respondent baseline update is unchanged from V9:
\begin{equation}
    \delta_i\mid\cdot\sim N(m_{\delta_i},V_{\delta_i}),
\end{equation}
where
\begin{align}
    V_{\delta_i}
    &=\left(P/\sigma_\varepsilon^2+1/\sigma_\delta^2\right)^{-1},\\
    m_{\delta_i}
    &=V_{\delta_i}\sigma_\varepsilon^{-2}
    \sum_{j=1}^P(x_{ij}-\lambda_i^\top\eta_j).
\end{align}

\paragraph{(7) $\lambda_i$.}
The loading row update is unchanged from V9:
\begin{equation}
    \lambda_i\mid\cdot\sim N_r(m_{\lambda_i},V_{\lambda_i}),
\end{equation}
where
\begin{align}
    V_{\lambda_i}
    &=\left(\tau_\lambda^{-2}I_r+
    \sigma_\varepsilon^{-2}\sum_{j=1}^P\eta_j\eta_j^\top\right)^{-1},\\
    m_{\lambda_i}
    &=V_{\lambda_i}\sigma_\varepsilon^{-2}
    \sum_{j=1}^P\eta_j(x_{ij}-\delta_i).
\end{align}

\paragraph{(8) $\sigma_\varepsilon^2$.}
The shared residual variance update is unchanged from V9:
\begin{equation}
    \sigma_\varepsilon^2\mid\cdot
    \sim
    IG\left(
    a_\varepsilon+\frac{nP}{2},
    b_\varepsilon+\frac{1}{2}
    \sum_{i=1}^n\sum_{j=1}^P
    (x_{ij}-\delta_i-\lambda_i^\top\eta_j)^2
    \right).
\end{equation}

\paragraph{(9) $\sigma_\delta^2$.}
\begin{equation}
    \sigma_\delta^2\mid\cdot
    \sim
    IG\left(a_\delta+\frac n2,
    b_\delta+\frac{1}{2}\sum_{i=1}^n\delta_i^2\right).
\end{equation}

\section{Split--merge move}
\label{sec:split-merge}

The split--merge move is a Metropolis--Hastings update on the partition
$c$, conditional on the current factor scores $\eta$.  It uses the collapsed
target \eqref{eq:log-collapsed-target}, where $\rho$ and
$(\mu_l,\Sigma_l)$ have been integrated out.

\subsection{Anchor selection}

Select two distinct item indices $u$ and $v$ uniformly at random from
$\{1,\dots,P\}$.  If $c_u=c_v$, attempt a split move.  If $c_u\neq c_v$,
attempt a merge move.  The anchor selection probability is symmetric and
cancels in the Metropolis--Hastings ratio.

\subsection{Split proposal}

Suppose $c_u=c_v=a$.  Let
\[
    A=\{j:c_j=a\}
\]
be the cluster to be split.  Choose an empty component label
$b\in\{l:n_l=0\}$ uniformly.  If no empty component exists, skip the split
proposal.

Initialize two temporary clusters by forcing the anchors to be separated:
\[
    c_u'=a,
    \qquad
    c_v'=b.
\]
For each item $j\in A\setminus\{u,v\}$, assign $j$ to either $a$ or $b$ by
a restricted NIW-collapsed Gibbs scan:
\begin{align}
    q_j(a)&\propto
    (n_{a,-j}^{\mathrm{temp}}+e_0)
    p(\eta_j\mid\eta_{a,-j}^{\mathrm{temp}}),\\
    q_j(b)&\propto
    (n_{b,-j}^{\mathrm{temp}}+e_0)
    p(\eta_j\mid\eta_{b,-j}^{\mathrm{temp}}),
\end{align}
where the predictive density is the NIW Student-$t$ density in
\eqref{eq:niw-predictive}, computed using the temporary cluster contents.
Let $p_j^{\mathrm{fwd}}$ denote the normalized probability of the assignment
actually selected for item $j$.  Then
\begin{equation}
    \log q(c'\mid c)
    =
    -\log E(c)+
    \sum_{j\in A\setminus\{u,v\}}
    \log p_j^{\mathrm{fwd}},
\end{equation}
where $E(c)=|\{l:n_l=0\}|$ is the number of empty labels before the split.
The reverse proposal is a deterministic merge, so
\[
    \log q(c\mid c')=0.
\]
The split is accepted with probability
\begin{equation}
    \alpha_{\mathrm{split}}
    =
    \min\left\{1,
    \exp\left[
    \log\pi(c'\mid\eta)-\log\pi(c\mid\eta)
    -\log q(c'\mid c)
    \right]
    \right\}.
\end{equation}

\subsection{Merge proposal}

Suppose $c_u=a$, $c_v=b$, and $a\neq b$.  Let
\[
    A=\{j:c_j=a\},
    \qquad
    B=\{j:c_j=b\}.
\]
The proposal merges $B$ into $A$:
\[
    c_j'=a,
    \qquad j\in A\cup B.
\]
Thus
\[
    \log q(c'\mid c)=0.
\]
The reverse proposal is a split of $A\cup B$ back into $A$ and $B$, using
anchors $u$ and $v$.  Its proposal probability is computed by running the
same restricted NIW-collapsed scan that would reproduce the original
assignment.  Let $p_j^{\mathrm{rev}}$ denote the normalized probability of
assigning item $j$ to its original side during this reverse split.  Then
\begin{equation}
    \log q(c\mid c')
    =
    -\log E(c')+
    \sum_{j\in (A\cup B)\setminus\{u,v\}}
    \log p_j^{\mathrm{rev}},
\end{equation}
where $E(c')$ is the number of empty labels after the merge.  The merge is
accepted with probability
\begin{equation}
    \alpha_{\mathrm{merge}}
    =
    \min\left\{1,
    \exp\left[
    \log\pi(c'\mid\eta)-\log\pi(c\mid\eta)
    +\log q(c\mid c')
    \right]
    \right\}.
\end{equation}

\section{One full MCMC sweep}

\begin{algorithm}[H]
\caption{One sweep of joint LSIRM + PPCA-mixture with NIW predictive and split--merge moves}
\begin{algorithmic}[1]
\State Update LSIRM non-position block.
\State Update $\{a_i^{(\ell)}\}$ via MH using the original LSIRM acceptance ratio.
\State Update $\{b_j^{(\ell)}\}$ via MH using the original LSIRM acceptance ratio.
\State Form $x_{ij}=\log\max\{\|a_i^{(\ell(j))}-b_j^{(\ell(j))}\|,\epsilon\}$ and row-center.
\If{iter $\geq T_{\mathrm{warm}}$}
    \State Update $\{\eta_j\}_{j=1}^P$ using \eqref{eq:eta-update-v10}.
    \State Update $\{c_j\}_{j=1}^P$ using the NIW-collapsed predictive update \eqref{eq:c-update-v10}.
    \For{$m=1,\dots,M_{\mathrm{SM}}$}
        \State Propose and accept/reject one split--merge move using \eqref{eq:log-collapsed-target}.
    \EndFor
    \State Update $\{\Sigma_l,\mu_l\}_{l=1}^{K^\star}$ using the NIW posterior draw \eqref{eq:mu-sigma-update-v10}.
    \State Optionally sample $\rho\mid c$ using \eqref{eq:rho-storage-v10} for storage.
    \For{$i=1,\dots,n$}
        \State Update $\delta_i$.
        \State Update $\lambda_i$.
        \State Accumulate $\mathrm{rss}\mathrel{+}=\sum_{j=1}^P(x_{ij}-\delta_i-\lambda_i^\top\eta_j)^2$.
    \EndFor
    \State Update $\sigma_\varepsilon^2$ using the accumulated residual sum of squares.
    \State Update $\sigma_\delta^2$.
    \State Update online co-cluster matrix $C_{jk}\leftarrow C_{jk}+\mathbb 1\{c_j=c_k\}$.
\EndIf
\end{algorithmic}
\end{algorithm}

\section{Expected effect of the V10 changes}

The V10 changes target a different failure mode from the V8-to-V9 change.
V9 prevents the loading matrix from collapsing by replacing
respondent-specific noise variances with a single PPCA residual variance.
V10 improves exploration of the clustering posterior once the factor scores
contain signal.

The NIW predictive update helps because it allocates item $j$ using
\[
    p(\eta_j\mid\eta_{\mathcal J_{l,-j}})
\]
rather than using a single sampled value of $(\mu_l,\Sigma_l)$.  Empty
components are assigned prior predictive mass and therefore can be
reoccupied.

The split--merge move helps because it proposes large changes to the
partition.  Single-site Gibbs must pass through low-probability intermediate
states to split a large cluster.  A split move can propose the two-cluster
configuration directly, and a merge move can remove spurious over-splitting.

Together, these changes are expected to reduce stickiness at $K_+=1$ or
$K_+=2$ and improve posterior exploration of $K_+\in\{3,4,\dots\}$.

\section{Cluster inference}

Cluster inference remains label-switching invariant.  The main posterior
summaries are:
\begin{itemize}[leftmargin=1.5em]
    \item Posterior similarity matrix
    $C_{jk}=\Pr(c_j=c_k\mid Y)$, computed online.
    \item Trace of occupied components
    $K_+^{(t)}=|\{l:n_l^{(t)}>0\}|$.
    \item A point partition obtained from the posterior similarity matrix,
    for example by hierarchical clustering on $1-C$ or by minimizing the
    posterior expected Binder loss.
\end{itemize}

\paragraph{Quality assessment.}  The point partition above should be
evaluated using \emph{$\eta$-space silhouette}: compute
$\mathrm{sil}_j$ for the chosen partition using Euclidean distance on the
posterior-mean factor scores $\hat\eta_j\in\mathbb R^r$.  Reporting
silhouette computed on $1-C$ instead is circular and tends to overstate
how well-separated clusters are.  See the verification checklist for
threshold conventions.

\section{Implementation notes}

The implementation requires the following changes to the V9 code.

\begin{enumerate}[leftmargin=2em]
    \item Replace the hyperparameter $V_0$ by $\kappa_0$ in
    \texttt{fmc\_hyper}.  Keep $m_0,\nu_0,S_0$.
    \item Add helper functions for NIW posterior statistics,
    multivariate Student-$t$ log density, NIW marginal log likelihood, and
    collapsed partition log target.
    \item Replace the current $c_j$ update based on
    $\rho_l\phi_r(\eta_j;\mu_l,\Sigma_l)$ by the NIW predictive update
    \eqref{eq:c-update-v10}.
    \item Replace the separate $\mu_l$ and $\Sigma_l$ updates by the NIW
    posterior draw \eqref{eq:mu-sigma-update-v10}.
    \item Move the $\rho$ update after the label updates, or keep it only as
    a monitoring draw.
    \item Add an optional argument such as \texttt{n\_split\_merge} or
    \texttt{sm\_prob} to control how often split--merge moves are attempted.
    \item Store split--merge acceptance counts separately for split and merge
    proposals.
\end{enumerate}

\section{Verification checklist}

After refitting V10, compare the following diagnostics with V9:
\begin{enumerate}[leftmargin=2em]
    \item $\|\Lambda\|_F$ trace and posterior mean: should remain non-degenerate.
    \item Reconstruction fraction of row-centered $X$ explained by
    $\Lambda\eta_j$: should remain meaningfully above zero.
    \item $K_+$ trace: should show more movement than V9, especially away
    from persistent $K_+=1$ or $K_+=2$.
    \item Split and merge acceptance rates: extremely low split acceptance
    suggests that the current factor scores support only coarse clusters or
    that the restricted Gibbs proposal is too weak.
    \item Co-cluster matrix: should show stable block structure if the
    posterior supports multi-cluster solutions.
    \item Empty-component reactivation: components should occasionally become
    occupied through the NIW prior predictive term, especially early in the
    chain.
    \item $S_0$ elicitation: the empirical within-cluster sd of $\eta_j$
    (computed using the modal labels) should be on the same order as
    $\sqrt{s_0/(\nu_0-r-1)}$.  If the prior implied sd is more than
    $\sim 5\times$ the empirical sd and split acceptance is low (e.g.\
    below $10\%$), tighten $S_0$ following \S\ref{sec:s0-elicitation}.
    \item Silhouette in factor-score space: compute the silhouette of the
    posterior partition using \emph{Euclidean distance between $\hat\eta_j$
    posterior means} (i.e.\ items live in $\mathbb R^r$, not in the
    co-cluster matrix).  Using $1-C$ as a dissimilarity, where $C$ is the
    posterior co-cluster matrix, is circular: $C$ is itself a smoothed
    function of how clustering on $\eta$ behaved across iterations, so
    the resulting "silhouette" mostly measures co-clustering consistency
    rather than separation in representation space.  Report (a) the
    overall mean silhouette, (b) the per-cluster silhouette mean, and
    (c) the fraction of items with negative silhouette.  Clusters with
    mean silhouette below $0.25$ should be flagged as weak; items with
    individually negative silhouette are candidates for re-assignment or
    indicate that the partition resolution is finer than the data
    supports.
    \item Co-clustering quality (two complementary measures):
    \begin{itemize}[leftmargin=1em]
        \item \textbf{Within-pair confidence ratio}
        \[
            r_{\rm wp}=\frac{\Pr(C_{jk}>0.8)}
                  {\text{frac.\ within-cluster pairs in final partition}}.
        \]
        Compares the fraction of pairs that the posterior puts together
        with high confidence to the fraction that the imposed partition
        puts together.  Values $r_{\rm wp}>0.7$ indicate that the
        partition is genuinely supported by the posterior; $r_{\rm wp}<0.3$
        indicates that the partition is largely a hierarchical-clustering
        artefact on a fuzzy $C$.
        \item \textbf{Posterior expected adjusted Rand}
        $\mathrm{PEAR}=E_{s,t}[\mathrm{ARI}(c^{(s)},c^{(t)})]$
        (Fritsch \& Ickstadt 2009): the average ARI between sampled
        partitions across MCMC iterations.  PEAR is label-switching-
        invariant: it asks whether the posterior is concentrated on one
        partition shape (high PEAR) or genuinely diffuse over partition
        space (low PEAR).  Threshold conventions: PEAR $>0.6$ = strong;
        $0.3$--$0.6$ = moderate; $<0.3$ = weak.

        \emph{When the two measures disagree.}  If $r_{\rm wp}$ is low but
        PEAR is high, individual pair co-cluster values are moderately
        ambiguous (e.g.\ $C_{jk}\!\approx\!0.4$--$0.6$) but the partition
        \emph{shape} is consistent across iterations.  This is meaningful
        and the partition can still be reported.  If both are low, the
        posterior is genuinely uncertain about the partition; report a
        $C$ heatmap and possibly only the most-stable subset of clusters
        that survives across iterations, rather than a single point
        partition.
    \end{itemize}
\end{enumerate}

\end{document}
